% GENERATED FILE. Edit canonical lessons, then run npm run generate:books.
% canonical-source-hash: fnv1a64:8fbe695b9d86227e
% canonical-lessons: BN-C16-na-verb, BN-C16-noy, BN-C16-nei, BN-C16-polar, BN-R16-three-noes

\chapter{Three Ways to Say No}
\label{ch:bn-three-ways-to-say-no}
\hlchaptermodality{27}

\begin{chapteropening}
\textbf{By the end of this chapter:} \emph{I can deny any sentence this book has taught me — one negator for a sentence with a verb, one for a sentence without, one for a thing that is not there — and I can turn every one of them into a yes-or-no question without moving a word.}

Three negators and one question shape, run against frames the book has been teaching since chapter three: the name frame denied with \bn{নয়}, the verb denied with \bn{না}, the missing thing denied with \bn{নেই}, and each of the three asked back with a rising voice.
\end{chapteropening}

\section[nā]{\bn{না} (nā) — the word that goes after}

\label{lesson:BN-C16-na-verb}



\begin{quote}
The first chapter of this book taught \textbf{\bn{না}} as an answer, and twenty-six chapters have gone by without it doing anything else. It is about to do the largest job in the language.

\begin{itemize}
\raggedright
  \item \emph{Recall:} two chapters back, the cloth word — \emph{kāpoṛ}
\end{itemize}
\end{quote}

\subsection*{You'll want to know: \bn{না} after the verb}
\begin{quote}
\textbf{\bn{আমি} \bn{জানি} \bn{না}} — \emph{āmi jāni nā} — \textbf{I do not know.}
\end{quote}

\begin{quote}
\textbf{\bn{আমি} \bn{ভাত} \bn{খাই} \bn{না}} — \emph{āmi bhāt khāi nā} — \textbf{I do not eat rice.}
\end{quote}

Take any sentence in this book, keep every word exactly where it stands, and put \textbf{\bn{না}} at the end. That is the whole operation. English drops \emph{not} in front of the verb and rebuilds the sentence around it — \emph{I know} becomes \emph{I do not know}, with a \emph{do} that came from nowhere. Bengali adds one syllable and moves nothing.

\begin{cousinweb}
\textbf{\bn{না}} is one of the oldest words a reader of English owns. Sanskrit \emph{na}, Latin \emph{ne-}, Greek \emph{nē-}, Old English \emph{ne}: it is the syllable behind \textbf{no}, \textbf{not}, \textbf{never}, \textbf{none} and \textbf{nay}. Chapter five already handed it over inside \textbf{\bn{কোনো} \bn{ব্যাপার} \bn{না}} (\emph{kono bæpār nā}), \textquotedblleft{}no matter\textquotedblright{} — a whole phrase ending in the negator, exactly where this lesson puts it.

One sentence this book cannot yet print is \emph{āmi bujhi nā}, \textquotedblleft{}I do not understand\textquotedblright{} — the most useful sentence a beginner can carry. The verb is taught, in chapter fourteen; the letter in the middle of it is not, so the words stay in romanization until a script lesson pays for that letter.
\end{cousinweb}

\subsection*{Your turn}
Say these aloud:

\begin{itemize}
\raggedright
  \item \emph{āmi jāni} — then again with the negator — \emph{āmi jāni nā}
  \item \emph{āmi bānglā bôli nā}
  \item \emph{āmi khāi nā} — I do not eat
  \item the sentence in romanization — \emph{āmi bujhi nā}
\end{itemize}

\begin{tcolorbox}[breakable,colback=teal!4,colframe=teal!35!black,title={Before you move on}]
Where in the sentence does \textbf{\bn{না}} stand? (\textbf{After the verb, at the end} — nothing else moves.) What did chapter one teach it as? (\textbf{The answer \textquotedblleft{}no\textquotedblright{}} — one word doing two jobs.) And what does English have to borrow to negate \emph{I know}? (\textbf{A verb \emph{do} that means nothing} — Bengali borrows nothing.)

Source: \href{https://www.unicode.org/charts/PDF/U0980.pdf}{Unicode Bengali chart}.
\end{tcolorbox}
\section[nôy]{\bn{নয়} (nôy) — denying a sentence with no verb in it}

\label{lesson:BN-C16-noy}



\begin{quote}
The lesson before this one put \textbf{\bn{না}} after a verb. Now try it on \textbf{\bn{আমার} \bn{নাম} \bn{রাম}} — and notice there is no verb there to put anything after.

\begin{itemize}
\raggedright
  \item \emph{Recall:} two chapters back, the shirt word — \emph{jāmā}
\end{itemize}
\end{quote}

\subsection*{You'll want to know: \bn{নয়}}
\begin{quote}
\textbf{\bn{আমার} \bn{নাম} \bn{রাম} \bn{নয়}} — \emph{āmār nām rām nôy} — \textbf{my name is not Ram.}
\end{quote}

\begin{quote}
\textbf{\bn{আমি} \bn{ভালো} \bn{নই}} — \emph{āmi bhālo nôi} — \textbf{I am not well.}
\end{quote}

Chapter three taught the frame \textbf{\bn{আমার} \bn{নাম} …} and said plainly that it carries no word for \emph{is}. Bengali leaves the present-tense \emph{is} out. To deny that sentence you cannot negate a verb that was never said, so the language supplies one, and it supplies it only here: \textbf{\bn{নয়}}.

\begin{grammarlens}[title={Grammar Lens: the form follows the person}]
\noindent
\begin{tabularx}{\linewidth}{@{}>{\raggedright\arraybackslash}X>{\raggedright\arraybackslash}X@{}}
\toprule
\textbf{about} & \textbf{say} \\
\midrule
me & \textbf{\bn{নই}} (\emph{nôi}) \\
you, or anybody else & \textbf{\bn{নয়}} (\emph{nôy}) \\
\bottomrule
\end{tabularx}

That is the whole paradigm a beginner needs. \textbf{\bn{আমি} \bn{ভালো} \bn{নই}} denies your own state; \textbf{\bn{রাম} \bn{ভালো} \bn{নয়}} denies somebody else's. Notice what the affirmative looks like beside it: \textbf{\bn{আমি} \bn{ভালো}} is a complete Bengali sentence with no verb at all. The copula in this language is a ghost — it is audible only when you refuse it.
\end{grammarlens}

\begin{sounds}
\textbf{\bn{নয়}} ends in \textbf{\bn{য়}}, which is \textbf{\bn{য}} wearing the dot chapter ten taught. Without the dot the letter says \emph{j}; with it, a glide, so the word is \emph{nôy} and not \emph{nôj}. The dot is doing exactly the job that lesson gave it.
\end{sounds}

\subsection*{Your turn}
Say these aloud:

\begin{itemize}
\raggedright
  \item \emph{āmār nām rām nôy}
  \item \emph{āmi bhālo nôi} — then the affirmative it denies, \emph{āmi bhālo}
  \item the verb kind beside the verbless kind — \emph{āmi jāni nā … āmi bhālo nôi}
\end{itemize}

\begin{tcolorbox}[breakable,colback=teal!4,colframe=teal!35!black,title={Before you move on}]
Which negator does a sentence with a verb take, and which does a sentence without one? (\textbf{\bn{না}} after the verb; \textbf{\bn{নয়}} when there is none.) Which form goes with \textbf{\bn{আমি}}? (\textbf{\bn{নই}}.) And where is the \emph{is} in \textbf{\bn{আমি} \bn{ভালো}}? (\textbf{Nowhere} — it exists only to be denied.)

Source: \href{https://www.unicode.org/charts/PDF/U0980.pdf}{Unicode Bengali chart}.
\end{tcolorbox}
\section[nei]{\bn{নেই} (nei) — the third negator, for things that are not there}

\label{lesson:BN-C16-nei}



\begin{quote}
Two negators so far: one after a verb, one for a sentence with no verb. A third is coming, and it exists because Bengali refuses to build this one out of the first two.

\begin{itemize}
\raggedright
  \item \emph{Recall:} two chapters back, the woven strip of cloth — \emph{shāṛi}
\end{itemize}
\end{quote}

\subsection*{You'll want to know: \bn{নেই}}
\begin{quote}
\textbf{\bn{আমার} \bn{চা} \bn{নেই}} — \emph{āmār chā nei} — \textbf{I have no tea.}
\end{quote}

\begin{quote}
\textbf{\bn{আমার} \bn{ভাই} \bn{নেই}} — \emph{āmār bhāi nei} — \textbf{I have no brother.}
\end{quote}

Chapter five taught \textbf{\bn{আমি} \bn{ভালো} \bn{আছি}}, \textquotedblleft{}I am well,\textquotedblright{} on the stem \textbf{\bn{আছ}-}, which means \emph{to be present}. \textbf{\bn{নেই}} is that stem's denial, and it comes as a single word rather than as \textbf{\bn{আছি} \bn{না}}, which no Bengali says.

\begin{grammarlens}[title={Grammar Lens: having, in a language with no verb for it}]
Bengali has no verb \emph{to have}. It says \emph{of-me tea is-present}, using the same \textbf{\bn{আমার}} chapter three taught for \emph{my}. So \textbf{\bn{আমার} \bn{চা} \bn{নেই}} is literally \emph{of me, tea is not present} — and the same shape covers \textbf{I have no tea} and \textbf{there is no tea of mine}, because in Bengali those were never different sentences.

\textbf{\bn{নেই}} does not change for the person. \textbf{\bn{আমার} \bn{নেই}}, \textbf{\bn{তোমার} \bn{নেই}}: one form for everybody, which makes it the easiest of the three negators and the one a traveller reaches for first.
\end{grammarlens}

\subsection*{Your turn}
Say these aloud:

\begin{itemize}
\raggedright
  \item \emph{āmār chā nei}
  \item \emph{āmār jôl nei} — I have no water
  \item all three in a row — \emph{jāni nā … bhālo nôi … chā nei}
\end{itemize}

\begin{tcolorbox}[breakable,colback=teal!4,colframe=teal!35!black,title={Before you move on}]
Which stem is \textbf{\bn{নেই}} the denial of? (\textbf{\bn{আছ}-}, the \emph{being present} stem behind \textbf{\bn{আছি}}.) Why is there no \textbf{\bn{আছি} \bn{না}}? (\textbf{Because the negator and the stem fused into one word} long ago.) And how does Bengali say \emph{I have}? (\textbf{With \bn{আমার}} — \emph{of me}, and no verb \emph{to have} at all.)

Source: \href{https://www.unicode.org/charts/PDF/U0980.pdf}{Unicode Bengali chart}.
\end{tcolorbox}
\section[tumi bhālo?]{\bn{তুমি} \bn{ভালো}? — a question with nothing added}

\label{lesson:BN-C16-polar}



\begin{quote}
Three chapters have taught you how to say \textbf{\bn{হ্যাঁ}} and \textbf{\bn{না}} to a question nobody in this book has been able to ask.

\begin{itemize}
\raggedright
  \item \emph{Recall:} two chapters back, the eye-thing — \emph{chôshmā}
\end{itemize}
\end{quote}

\subsection*{You'll want to know: the rising voice}
\begin{quote}
\textbf{\bn{তুমি} \bn{ভালো}} — \emph{tumi bhālo.} — \textbf{you are well.}
\end{quote}

\begin{quote}
\textbf{\bn{তুমি} \bn{ভালো}?} — \emph{tumi bhālo?} — \textbf{are you well?}
\end{quote}

Those are the same four syllables in the same order. The only difference is that the second one ends higher than it began. English rebuilds the sentence — \emph{you are} becomes \emph{are you} — and Bengali does not touch it.

\begin{grammarlens}[title={Grammar Lens: what the rise turns into a question}]
Every frame in this book can now be asked as well as said, at no cost:

\noindent
\begin{tabularx}{\linewidth}{@{}>{\raggedright\arraybackslash}X>{\raggedright\arraybackslash}X@{}}
\toprule
\textbf{said} & \textbf{asked} \\
\midrule
\textbf{\bn{আমার} \bn{চা} \bn{নেই}} & \textbf{\bn{তোমার} \bn{চা} \bn{নেই}?} \\
\textbf{\bn{তুমি} \bn{ভাত} \bn{খাও}} & \textbf{\bn{তুমি} \bn{ভাত} \bn{খাও}?} \\
\bottomrule
\end{tabularx}

A negative asked this way is the ordinary way to check something you suspect: \textbf{\bn{তোমার} \bn{চা} \bn{নেই}?} — \emph{you have no tea?} And the answer is one of the two words chapter one taught.

Bengali also has a written question particle, \textbf{\bn{কি}}, which chapter three taught as the word for \emph{what}. In a yes-or-no question it does the job the rising voice does in speech. The voice is what a speaker needs first.
\end{grammarlens}

\subsection*{Your turn}
Say these aloud:

\begin{itemize}
\raggedright
  \item \emph{tumi bhālo.} flat — then \emph{tumi bhālo?} rising
  \item \emph{tomār chā nei?} — and answer it \emph{hyã} or \emph{nā}
  \item \emph{āpni bhālo?} — the respectful one
\end{itemize}

\begin{tcolorbox}[breakable,colback=teal!4,colframe=teal!35!black,title={Before you move on}]
What moves when a Bengali statement becomes a yes-or-no question? (\textbf{Nothing} — only the pitch at the end.) What does English have to move? (\textbf{The verb, in front of the subject}.) And which two words answer it? (\textbf{\bn{হ্যাঁ} and \bn{না}}, from the first chapter.)

Source: \href{https://www.unicode.org/charts/PDF/U0980.pdf}{Unicode Bengali chart}.
\end{tcolorbox}
\section[nā · nôy · nei]{Three ways to say no, and one way to ask}

\label{lesson:BN-R16-three-noes}



\begin{quote}
Three negators arrived in this chapter and each one owns a different kind of sentence. Name them in the order you met them, and name what each one denies.
\end{quote}

\subsection*{Your turn}
\begin{itemize}
\raggedright
  \item \emph{Say it:} the sentence with a verb — \emph{āmi jāni nā}
  \item \emph{Say it:} the sentence with none — \emph{āmār nām rām nôy}
  \item \emph{Say it:} the thing that is missing — \emph{āmār chā nei}
  \item \emph{Say it:} each of those three again with a rising end, and hear a question
  \item \emph{Recall:} the sentence chapter eleven taught, then deny it — \emph{āmi bānglā bôli nā}
\end{itemize}

\begin{tcolorbox}[breakable,colback=teal!4,colframe=teal!35!black,title={Before you move on}]
Which negator would you use to deny \textbf{\bn{আমার} \bn{নাম} \bn{রাম}}, which for \textbf{\bn{আমি} \bn{জানি}}, and which for \textbf{\bn{আমার} \bn{চা} \bn{আছে}}? (\textbf{\bn{নয়}}, then \textbf{\bn{না}}, then \textbf{\bn{নেই}}.) What has to change to turn any of the three into a question? (\textbf{The pitch at the end, and nothing else}.)
\end{tcolorbox}
